% 第12章（本部分第5章）图：分块消元。增广分块矩阵 [B C | f; D E | g] 的第二块行减去 D B^{-1} 乘第一块行，
% 左下块化为零，右下块成为 Schur 补 S=E-DB^{-1}C，右端成为 h=g-DB^{-1}f。块尺寸示意 p>q。
\begin{tikzpicture}[x=1mm,y=1mm,>=Stealth,
  blk/.style={rectangle,draw=FigureStroke,fill=white,line width=.6pt,inner sep=0pt,font=\figurefont\fontsize{8}{10}\selectfont},
  note/.style={font=\FigureNoteFont\fontsize{7.5}{9}\selectfont,text=FigureMuted}]
% block geometry: p=16mm, q=8mm, rhs width 5mm
\newcommand{\augmat}[9]{%
 % #1 x-origin, #2..#9: labels/styles  B C f D E g + fill for D-block, fill for E-block
 \node[blk,minimum width=16mm,minimum height=16mm] at (#1+8,-8) {$\bm B$};
 \node[blk,minimum width=8mm,minimum height=16mm] at (#1+20,-8) {$\bm C$};
 \node[blk,minimum width=5mm,minimum height=16mm] at (#1+28.5,-8) {$\bm f$};
 \node[blk,minimum width=16mm,minimum height=8mm,#8] at (#1+8,-20) {#2};
 \node[blk,minimum width=8mm,minimum height=8mm,#9] at (#1+20,-20) {#3};
 \node[blk,minimum width=5mm,minimum height=8mm,#9] at (#1+28.5,-20) {#4};
}
\augmat{0}{$\bm D$}{$\bm E$}{$\bm g$}{}{}{}{}{}
\augmat{62}{$\bm0$}{$\bm S$}{$\bm h$}{}{}{}{fill=FigureGreenFill,draw=FigureGreen}{fill=FigureYellowFill,draw=FigureYellow}
% size marks
\node[note,anchor=east] at (-1,-8) {$p$};
\node[note,anchor=east] at (-1,-20) {$q$};
% arrow
\draw[->,FigureStroke,line width=.8pt] (35,-14)--(59,-14);
\node[note,align=center,anchor=south] at (47,-13) {第二块行减去\\$\bm D\bm B^{-1}$乘第一块行};
\node[note,align=center,anchor=north] at (47,-15.5) {用$\bm B$的因子\\求$\bm Z,\bm z$};
% formulas
\node[note,anchor=north,align=left] at (77.5,-26.5) {$\bm S=\bm E-\bm D\bm Z,\ \bm Z=\bm B^{-1}\bm C$\\$\bm h=\bm g-\bm D\bm z,\ \bm z=\bm B^{-1}\bm f$};
\end{tikzpicture}
